\documentclass[11pt,a4paper]{article}
\usepackage[T1]{fontenc}
\usepackage[utf8]{inputenc}
\usepackage{lmodern}
\usepackage{amsmath,amssymb}
\usepackage[margin=25mm]{geometry}
\usepackage[hidelinks]{hyperref}
\hypersetup{pdftitle={Lifting a Lipschitz Newton certificate to quantum mechanics},pdfauthor={navstokgap research working notes}}
\newcommand{\tightlist}{\setlength{\itemsep}{0pt}\setlength{\parskip}{0pt}}
\setlength{\emergencystretch}{3em}
\setcounter{secnumdepth}{0}
\begin{document}
\hypertarget{lifting-a-lipschitz-newton-certificate-to-quantum-mechanics}{%
\section{Lifting a Lipschitz Newton certificate to quantum
mechanics}\label{lifting-a-lipschitz-newton-certificate-to-quantum-mechanics}}

\textbf{Result, 2026-10-27; written derivation, unrefereed.} A classical
kick--drift existence certificate for a one-dimensional conservative
force with a global Lipschitz constant supplies the model-dependent
estimates for a specified noncommutative quantum construction. For every
supplied \(\hbar>0\) it gives the Schrödinger representation, an
essentially self-adjoint Hamiltonian, a strongly continuous unitary
group, and arbitrary-partition convergence of the quantum kick--drift
products. On coherent preparations the same classical polygons and
quantum products agree at their vertices to \(O(\sqrt\hbar)\), uniformly
in the mesh on a fixed time interval. Together with the classical
certificate this gives the requested joint observable limit.

This is a \textbf{bounded instance}, not a proof of the universal
\href{classical-quantum-existence-conjecture.md}{constructive lifting
conjecture}. The input includes the conservative action and a globally
Lipschitz force; the background analytic toolkit includes Nelson's
commutator theorem and the graph-domain propagation result cited below.
No lower bound on \(V\) or derivatives beyond the Lipschitz force are
assumed. This real-time theorem does not construct a Euclidean
probability measure for a potential unbounded below.

The loss of the earlier \(O(\hbar)\) rate is necessary at this level of
regularity: \(V(q)=q|q|/2\) has globally Lipschitz force, but a centered
coherent packet has a momentum error of exact order \(\sqrt\hbar\) at
small fixed positive times. No claim of literature priority is made.

\hypertarget{the-input-certificate-and-quantum-target}{%
\subsection{1. The input certificate and quantum
target}\label{the-input-certificate-and-quantum-target}}

Let \(m>0\), \(V\in C^1(\mathbb R;\mathbb R)\), and suppose
\[|V'(x)-V'(y)|\le L|x-y|,\qquad L\ge0.\tag{1}\] Thus \(V\in C^{1,1}\)
and \(V''\) exists almost everywhere with \(|V''|\le L\). The input is
the mechanical action with kinetic term \(p^2/(2m)\) and potential
\(V\), the initial point \(z_0=(q_0,p_0)\), and the classical
kick--drift polygons \[
 p_{j+1}^{\rm cl}=p_j^{\rm cl}-s_jV'(q_j^{\rm cl}),\qquad
 q_{j+1}^{\rm cl}=q_j^{\rm cl}+s_jp_{j+1}^{\rm cl}/m.\tag{2}
\] Take their linear interpolation for a partition \(\pi\) of \([0,T]\).
The usual Lipschitz certificate supplies a unique limit \(z(t)\) and
\[\sup_{t\le T}\|z_\pi^{\rm cl}(t)-z(t)\|_\omega
 \le R_T^{\rm cl}|\pi|,\qquad
 \|(q,p)\|_\omega=|q|+\frac{|p|}{m\omega}.\tag{3}\] For example, the
proof in the \href{balaban-ode-lesson.md}{ODE lesson} uses a common
bounded region, force-balance residual and Gronwall stability. Here
\(\omega>0\) is a fixed auxiliary unit; \(R_T^{\rm cl}\) depends only on
the classical data and \(T\). No quantum convergence is included in the
certificate.

For each \(\hbar>0\) set, on \(L^2(\mathbb R)\), \[
 \begin{gathered}
 Qf(q)=qf(q),\quad P=-i\hbar\partial_q,\quad
 H_\hbar=P^2/(2m)+V(Q),\\
 K_{\hbar,s}=e^{-isV(Q)/\hbar},\quad
 D_{\hbar,s}=e^{-isP^2/(2m\hbar)}.
 \end{gathered}\tag{4}
\] The self-adjoint \(Q,P\) obey the Weyl relations and
\([Q,P]=i\hbar I\) on Schwartz space. Unlike the Koopman construction,
these observables are noncommutative.

\hypertarget{the-same-force-bound-controls-the-operator-domains}{%
\subsection{2. The same force bound controls the operator
domains}\label{the-same-force-bound-controls-the-operator-domains}}

Fix \(\hbar>0\) and calibrated units and put \(N=1+Q^2+P^2\). The input
(1) implies \[
 |V'(q)|\le |V'(0)|+L|q|,\qquad
 |V(q)|\le |V(0)|+|V'(0)||q|+\tfrac L2q^2.\tag{5}
\] Using the harmonic-oscillator graph bounds for \(Q^2,P^2,QP,PQ\),
these estimates and the almost-everywhere bound on \(V''\) imply on
\(\mathcal S(\mathbb R)\) \[
 \|H_\hbar f\|\le C_\hbar\|Nf\|,\qquad
 \|[H_\hbar,N]f\|\le C_\hbar\|Nf\|.\tag{6}
\] The commutator is a weak-derivative identity: \[
 [H_\hbar,N]
 =-\frac{i\hbar}{m}(PQ+QP)
    +i\hbar\big(V'(Q)P+PV'(Q)\big).\tag{7}
\] Products with \(V'\) have the weak derivative \(V''\); no third
derivative is used. For \(f\in\mathcal S\), \(H_\hbar f\in D(N)\), so
(6) gives \[
 |\langle H_\hbar f,N^2f\rangle-\langle N^2f,H_\hbar f\rangle|
 \le2C_\hbar\langle f,N^2f\rangle.\tag{8}
\] Schwartz space is a core for \(N^2\). Nelson's commutator theorem
therefore makes \(H_\hbar\) essentially self-adjoint on \(\mathcal S\).
With \(G=N^2-\inf\sigma(N^2)\), the same estimate and
\href{https://link.springer.com/article/10.1007/s00220-025-05282-w}{Energy-Limited
Quantum Dynamics, Theorems 3.7--3.8} (statements and hypotheses read)
show that \(U_\hbar(t)=e^{-itH_\hbar/\hbar}\) preserves \(D(N)\) and is
strongly continuous in its \(N\)-graph norm. The constants in (6)--(8)
follow from \(m,\hbar,L,V(0),V'(0)\); they are not additional
assumptions on the mechanical model.

\hypertarget{a-fixed-action-quantum-product-estimate}{%
\subsection{3. A fixed-action quantum product
estimate}\label{a-fixed-action-quantum-product-estimate}}

The force certificate also controls the commutator of the two factors.
Write \(T_\hbar=P^2/(2m)\) and, initially on Schwartz space,
\[[T_\hbar,V(Q)]
 =-\frac{i\hbar}{2m}\big(PV'(Q)+V'(Q)P\big).\tag{9}\] It extends as a
bounded map \(D(N)\to L^2\). To justify the product argument without
assuming higher derivatives, mollify \(V\) with an even compactly
supported unit-mass mollifier. The smooth \(V_n\) satisfy
\(\|V_n''\|_\infty\le L\) and have uniformly linearly growing first
derivatives. Moreover, Taylor's inequality gives \[\|V_n-V\|_\infty
 \le \tfrac12L\,\varepsilon_n^2
          \int y^2\rho(y)\,dy\longrightarrow0.\tag{10}\] This uniform
bound holds even when \(V\) is unbounded below.

For these \(V_n\), define the finite constants \[
 \eta_T=\hbar^{-1}\|[T_\hbar,N]N^{-1}\|,\quad
 \eta_V=\sup_n\hbar^{-1}\|[V_n(Q),N]N^{-1}\|,\quad
 B=\sup_n\|[T_\hbar,V_n(Q)]N^{-1}\|.\tag{11}
\] They are finite by (5) and the graph bounds above; their values can
be bounded from the classical data \(m,L,V'(0)\) and the chosen
\(\hbar\) and units. The \(N\)-graph norms of the kinetic, potential and
full groups are bounded by \(e^{\eta_T|t|}\), \(e^{\eta_V|t|}\), and
\(e^{(\eta_T+\eta_V)|t|}\), respectively.

For smooth \(V_n\), differentiation of \(U_n(s-r)D_rK_{n,r}\) and a
second Duhamel integration yield, on Schwartz vectors, \[
 \|(D_sK_{n,s}-U_n(s))f\|
 \le\frac{Bs^2}{2\hbar^2}
       e^{(\eta_T+\eta_V)s}\|Nf\|.\tag{12}
\] The left factors in the integral are unitary; the commutator (9) acts
on \(D_uK_{n,r}f\). Density extends (12) to \(D(N)\). Telescoping along
the exact \(U_n\) orbit and using the graph bound gives \[
 \|U_{n,\pi}(T)f-U_n(T)f\|
 \le\frac{BT\,e^{(\eta_T+\eta_V)T}}{2\hbar^2}
      |\pi|\,\|Nf\|,\qquad f\in D(N).\tag{13}
\] Equation (10) makes \(U_n(t)\to U_\hbar(t)\) and each finite product
\(U_{n,\pi}(T)\to U_{\hbar,\pi}(T)\) in operator norm for fixed
\(\hbar,T,\pi\), by the bounded-perturbation Duhamel estimate. The
constants in (13) are uniform in \(n\), so (13) passes to \(V\). It
proves arbitrary-partition strong convergence for \(f\in D(N)\); density
and unitarity extend convergence to all \(L^2\) vectors. In particular,
the conjecture note's fixed-\(\hbar\) Cauchy criterion holds, with the
explicit graph-domain rate in (13). The factor \(\hbar^{-2}B\) need not
remain bounded as \(\hbar\downarrow0\).

\hypertarget{the-classical-and-quantum-polygons}{%
\subsection{4. The classical and quantum
polygons}\label{the-classical-and-quantum-polygons}}

Use the coherent packet of the
\href{quantum-blocks-construct-newton.md}{quantum-blocks note}, equation
(16), and let \((q_j,p_j)\) be the means after the first \(j\) quantum
cells. The exact identities \[
 K_s^*PK_s=P-sV'(Q),\qquad D_s^*QD_s=Q+sP/m\tag{14}
\] hold on the propagated states; the \(C^{1,1}\) potential preserves
the needed graph domain. Define \[
 W_j=\operatorname{sd}Q_j+\frac{\operatorname{sd}P_j}{m\omega},
 \quad C=\sqrt{\frac2{m\omega}},\quad
 \beta=\omega+\frac{L}{m\omega}.\tag{15}
\] As in the earlier note, (1) and (14) give
\[W_j\le C\sqrt\hbar\,e^{\beta t_j}.\tag{16}\] No higher potential
derivative or assumption of Gaussian shape after the first step is used.

Put \(e_j=\|(q_j,p_j)-(q_j^{\rm cl},p_j^{\rm cl})\|_\omega\). For each
kick,
\(|\langle V'(Q)\rangle-V'(q_j^{\rm cl})|  \le L(\operatorname{sd}Q_j+|q_j-q_j^{\rm cl}|)\).
The subsequent drift yields \[
 e_{j+1}\le
 (1+s_j\omega)\!\left(1+\frac{s_jL}{m\omega}\right)e_j
 +\frac{s_jL}{m\omega}(1+s_j\omega)W_j.\tag{17}
\] Since \(e_0=0\), (16) and \(1+x\le e^x\) imply the
\textbf{mesh-uniform} bound \[
 e_j\le \frac{LC}{m\omega}\,t_j e^{\beta t_j}\sqrt\hbar.
 \tag{18}
\] Combining this with the classical certificate (3) gives \[
 \|(q_j,p_j)-z(t_j)\|_\omega
 \le \frac{LC}{m\omega}T e^{\beta T}\sqrt\hbar
       +R_T^{\rm cl}|\pi|.\tag{19}
\] For \(b\in C_c^\infty(\mathbb R^2)\), the Fourier/Weyl argument in
\href{groupoid-quantum-newton-states.md}{the groupoid note}, before its
second-order refinement, bounds the expectation error by a
symbol-dependent constant times the sum of \(W_j\) and (19). Therefore
the observable limit required by the conjecture holds at all partition
vertices and along every joint schedule \(\hbar,|\pi|\to0\).

The proof uses the classical \textbf{input bound} (1) for both the graph
and packet estimates, while the certificate's polygon error (3) controls
recovery of the classical flow. The graph estimate (13) is a separate
fixed-\(\hbar\) quantum estimate derived using the declared operator
toolkit; it is not assumed in (3).

\hypertarget{why-the-square-root-rate-cannot-generally-be-improved}{%
\subsection{5. Why the square-root rate cannot generally be
improved}\label{why-the-square-root-rate-cannot-generally-be-improved}}

Set \(m=1\) and \[V(q)=\tfrac12q|q|,\qquad V'(q)=|q|.\tag{20}\] The
force is globally Lipschitz, and the classical solution from \((0,0)\)
is uniquely stationary. Define the unitary dilation
\((D_\hbar f)(q)=\hbar^{-1/4}f(q/\sqrt\hbar)\). Homogeneity gives
\[D_\hbar^{-1}H_\hbar D_\hbar=\hbar H_1,\qquad
 D_\hbar^{-1}P_\hbar D_\hbar=\sqrt\hbar\,P_1.\tag{21}\] For the centered
Gaussian \(D_\hbar g_1\), the exact quantum momentum is therefore
\(\langle P_\hbar\rangle_t=\sqrt\hbar\,f(t)\), where
\(f(t)=\langle P_1\rangle_t\). Ehrenfest's identity gives
\[f'(0)=-\langle |Q|\rangle_{g_1}=-1/\sqrt\pi.\tag{22}\] Thus
\(f(t)\ne0\) for every sufficiently small fixed \(t>0\). The mean error
is exactly a nonzero multiple of \(\sqrt\hbar\) there. This
counterexample does not contradict the \(O(\hbar)\) theorem for bounded
\(V'''\): \(V'\) here has a corner at the origin.

\hypertarget{consequence-for-state}{%
\subsection{6. Consequence for STATE}\label{consequence-for-state}}

The bounded certificate-to-quantum test is affirmative for the specified
one-dimensional conservative \(C^{1,1}\) action class. It constructs the
noncommutative target and a quantum product estimate from the classical
force data and the named analytic toolkit, and uses the classical
polygon certificate in the observable limit. The universal conjecture,
singular classical actions, quantum completion uniqueness and
independent positive-action necessity remain open.
\end{document}
